\chapter[Referêncial Teórico]{Referêncial Teórico}
\label{cap_exemplos}

\subsection{Fundamentos da Experimentação no Ensino de Ciências}
A experimentação laboratorial ocupa um papel central e inalienável na física como disciplina científica e no desenvolvimento individual de competências profissionais na educação em ciências \cite{blind}. O ensino de ciências, e em particular o de Física, é historicamente pautado no equilíbrio entre a formalização matemática e a experimentação prática, sendo a atividade experimental fundamental para converter abstrações em compreensão empírica. Longe de ser um mero complemento às aulas expositivas, o laboratório é o ambiente onde o estudante pode desenvolver uma visão física do mundo ao coletar, analisar e interpretar dados oriundos de suas próprias medições \cite{transforming}.

Sob a lente da epistemologia genética de Jean Piaget, a experimentação é proposta como uma atividade desequilibrante \cite{optics}. Ao manipular variáveis e observar fenômenos reais, o aluno é confrontado com situações que desafiam suas estruturas cognitivas prévias, forçando-o a rever e modificar esquemas mentais já construídos para acomodar a nova realidade observada. Complementarmente, na perspectiva vigotskiana, as atividades experimentais desencadeiam interações sociais oportunas entre alunos e professores. Essas trocas, ocorridas dentro da Zona de Desenvolvimento Proximal (ZDP), permitem que o aprendiz, mediado por um parceiro mais experiente na cultura científica, construa significados culturais e aproprie-se da linguagem e dos métodos da ciência.

A literatura destaca que o aprendizado de longo prazo é significativamente potencializado pela experimentação ativa direta \cite{foam}. No entanto, o modelo tradicional de "laboratório de receita de bolo" (cookbook labs), focado apenas na conformidade procedimental, tem sido criticado por falhar na promoção da aprendizagem conceitual profunda \cite{smartphone}. Em contrapartida, abordagens baseadas em projetos e investigações abertas incentivam os estudantes a agirem como cientistas: formulando perguntas, desenhando experimentos e lidando com as incertezas inerentes ao mundo real \cite{vidal}.

A construção do conhecimento físico no laboratório exige que o aluno consiga reconciliar a discrepância entre os modelos ideais dos livros didáticos e os dados experimentais reais, sujeitos a ruídos e variáveis parasitas \cite{kowalski}. Enquanto simuladores virtuais operam sob simplificações matemáticas, o contato com aparatos físicos — mesmo aqueles baseados em "toy models" ou kits de baixo custo — permite ao estudante experimentar a teoria da propagação de erros e a validação de modelos em condições não ideais. Conforme argumentado por Kowalski (2024), a compreensão de que as leis científicas são descrições aproximadas da realidade, constantemente refinadas mas nunca absolutas, é crucial para o desenvolvimento da alfabetização científica \cite{taylor}.

Além do ganho conceitual, as práticas laboratoriais visam ao desenvolvimento das chamadas Science Process Skills (SPS), que incluem a capacidade de observação, tomada de decisão, trabalho em equipe e comunicação de resultados. O uso de novas tecnologias, como microcontroladores e telemetria, não apenas motiva os alunos através da atratividade tecnológica, mas também os prepara para as demandas de uma força de trabalho qualificada e tecnologicamente letrada \cite{thermoelectric}. Assim, a experimentação real consolida-se como o elo necessário entre o pensamento formal e a prática científica autêntica.

\subsection{Evolução das Tecnologias de Laboratório: Do Local ao Remoto}
A infraestrutura laboratorial voltada ao ensino de ciências passou por transformações profundas nas últimas décadas, impulsionadas pela onipresença das Tecnologias Digitais de Informação e Comunicação (TDIC). Historicamente, o modelo predominante é o do acesso local com recursos reais, no qual estudantes e professores devem estar fisicamente presentes em um espaço comum para manipular aparatos e instrumentos. Contudo, esse modelo enfrenta limitações críticas de escalabilidade e espaço físico. Tais barreiras estruturais frequentemente resultam em laboratórios subequipados ou com hardware ocioso durante a maior parte do tempo letivo.

Para compreender a transição tecnológica, a literatura classifica os ambientes de experimentação em quatro categorias fundamentais: (a) Local access-real resource, o laboratório tradicional; (b) Local access-simulated resource, comum em cursos de computação; (c) Remote access-real resource, onde os experimentos reais são controlados à distância; e (d) Remote access-simulated resource, que abrange simuladores virtuais acessados via web \cite{transforming}. Embora os simuladores virtuais (SV) permitam uma visualização clara de fenômenos, eles operam sob simplificações matemáticas de modelos ideais, o que pode privar o aluno do contato com a incerteza e o ruído inerentes à natureza.

Neste contexto, os Laboratórios de Acesso Remoto (LAR) ou Remote Controlled Laboratories (RCL) surgiram como uma alternativa superior para a democratização do ensino experimental. Iniciadas de forma mais expressiva na década de 90, as primeiras propostas buscavam permitir que estudantes interagissem com equipamentos caros ou perigosos de qualquer lugar, eliminando barreiras geográficas e temporais \cite{termometria}. Sistemas pioneiros utilizaram arquiteturas baseadas em LabVIEW ou MATLAB Web Server (MWS) para mediar a comunicação entre o cliente e o hardware de coleta no servidor \cite{quadrant}. Essas tecnologias permitiram a criação de interfaces onde o usuário define parâmetros, aciona o experimento e recebe resultados em tempo real via navegador, sem a necessidade de softwares adicionais.

As vantagens documentadas dos LAR incluem a disponibilidade 24/7, o compartilhamento de recursos entre diferentes instituições e a segurança no manuseio de componentes sensíveis \cite{transforming}. Entretanto, as primeiras infraestruturas eram frequentemente monolíticas e "caixas-pretas", com experimentos codificados de forma rígida, o que dificultava a integração de novos aparatos e a manutenção por profissionais que não fossem especialistas em engenharia de software \cite{FREE}.

A evolução mais recente das tecnologias de laboratório foi dramaticamente acelerada pela pandemia de COVID-19 \cite{inverso}. O fechamento repentino dos campi universitários forçou instituições globais a migrarem para modelos de ensino híbrido e laboratórios distribuídos. Projetos como os CREATE labs demonstraram que a transformação de laboratórios tradicionais em experiências remotas-híbridas envolve não apenas o hardware, mas uma completa revisão pedagógica, utilizando ferramentas de colaboração online e câmeras em primeira pessoa para manter o engajamento do aluno \cite{create}.

Atualmente, a fronteira tecnológica foca em frameworks abertos e orquestrados, como o FREE (Framework for Remote Experiments in Education). Diferente dos sistemas isolados do passado, essas arquiteturas modernas defendem a separação clara entre a interface do usuário (UI), a lógica do servidor e o controle do aparato físico via protocolos leves como REST, HTTP e JSON. Essa mudança de paradigma permite que múltiplos nós sensores, baseados em microcontroladores de baixo custo como o ESP32, atuem de forma integrada em uma rede IoT, centralizando a telemetria em servidores locais (como o Raspberry Pi) que, por sua vez, podem ser integrados a sistemas como um todo. Assim, a evolução tecnológica caminha para a consolidação de laboratórios que não são apenas "remotos", mas orquestrados, escaláveis e pedagógica e tecnicamente integrados ao fluxo de trabalho digital da educação moderna.

\subsection{Microcontroladores e a Internet das Coisas (IoT) na Educação}
A introdução de plataformas de prototipagem eletrônica no ambiente educacional marcou uma transição tecnológica de instrumentos de medição isolados para nós inteligentes em uma rede de dados orquestrada. No centro dessa mudança estão os microcontroladores, que permitem que aparatos físicos sejam integrados a Sistemas de Informação, possibilitando a coleta e transmissão de telemetria de forma autônoma.

A utilidade dessas tecnologias no ensino de ciências transcende o uso simplificado por entusiastas, apresentando uma robustez técnica que desafia sistemas comerciais. Nino et al. (2014) demonstram que, ao otimizar o software em nível de registradores e interrupções, um microcontrolador de baixo custo pode estabilizar sistemas de alta complexidade com uma precisão de 200 kHz, superando em uma ordem de grandeza as especificações de fábrica de equipamentos proprietários. Essa capacidade de processamento permite que o hardware atue na digitalização de fenômenos físicos dinâmicos sem a necessidade de um computador de alto desempenho dedicado para cada experimento \cite{nino}.

Dentro do paradigma da IoT, o ESP32 destaca-se como uma plataforma versátil devido à sua conectividade sem fio (Wi-Fi e Bluetooth) nativa \cite{esp32}. Em arquiteturas laboratoriais modernas, esses dispositivos atuam na camada de borda (edge), sendo responsáveis pela aquisição de sinais analógicos e pela comunicação via protocolos leves como MQTT ou HTTP. Ao contrário de abordagens que utilizam o microcontrolador apenas para coleta serial, o modelo orquestrado delega a responsabilidade de controle e persistência a um servidor central (como o Raspberry Pi), permitindo que o laboratório opere como uma rede integrada de dispositivos conectados.

Essa integração permite que o docente gerencie o progresso de múltiplas bancadas simultaneamente através de interfaces web, eliminando o emaranhado de cabos e a ociosidade técnica de hardware que opera sem conectividade
. Assim, os microcontroladores deixam de ser apenas ferramentas de medição e passam a ser os alicerces de uma infraestrutura de informação escalável, essencial para o suporte a laboratórios híbridos e remotos.

\subsection{Arquiteturas de Orquestração e Gestão de Dados Centralizada}
A evolução dos laboratórios de ensino, de unidades autônomas para redes integradas de telemetria, exige uma mudança de paradigma na forma como os dados são capturados, processados e apresentados. Em vez de ferramentas de medição isoladas, as infraestruturas modernas operam sob o conceito de orquestração de dispositivos, onde múltiplos nós sensores cooperam dentro de uma hierarquia de rede para fornecer uma visão holística do ambiente experimental. Esta centralização é vital para superar as limitações de sistemas que dependem de conexões físicas individuais, permitindo que o docente gerencie o progresso de diversas bancadas simultaneamente através de um único ponto de controle.

Um pilar fundamental desta orquestração é o princípio do desacoplamento entre a Interface do Usuário (UI) e o controle do aparato físico. Segundo o framework FREE (Framework for Remote Experiments in Education), sistemas monolíticos e rigidamente codificados dificultam a manutenção e a escalabilidade do laboratório. A arquitetura orquestrada propõe uma separação clara onde a lógica de visualização no Dashboard Web pode ser alterada ou atualizada sem a necessidade de modificar o firmware dos microcontroladores de borda. Este modelo de computação em borda (edge computing) permite que o microcontrolador (ESP32) foque exclusivamente na aquisição de sinais em tempo real, enquanto o servidor central (Raspberry Pi) atua como um proxy agnóstico que garante a interoperabilidade entre diferentes linguagens e protocolos.

A segurança e a viabilidade técnica da rede também são fortalecidas pela gestão centralizada. Plataformas como a CEyeClon demonstram que arquiteturas centralizadas oferecem níveis de segurança superiores a aplicações web padrão por evitarem a necessidade de configurar IPs públicos, VPNs ou aberturas complexas de portas de firewall para cada experimento individual. Nestes frameworks modernos, os nós de borda iniciam a comunicação com o servidor, o que simplifica a infraestrutura de rede institucional e reduz a sobrecarga administrativa, mantendo a integridade da telemetria \cite{spring}.

Além do monitoramento em tempo real, a orquestração centralizada é a única forma de garantir a persistência de dados e a análise pedagógica posterior. Sem um servidor central dotado de banco de dados (como MySQL ou SQL Server), os dados gerados durante um fenômeno físico — como a descarga de um capacitor — seriam perdidos após a conclusão do ensaio. A centralização permite que o sistema armazene o histórico de múltiplas execuções, permitindo que alunos acessem resultados passados para revisão e que professores avaliem o desempenho técnico e conceitual de toda a turma de forma síncrona ou assíncrona.

Por fim, a fronteira definitiva da gestão centralizada reside na integração com o ecossistema educacional digital, especificamente com Sistemas de Gerenciamento de Aprendizagem (LMS). Através do protocolo LTI (Learning Tools Interoperability), a arquitetura orquestrada permite que o laboratório "converse" diretamente com plataformas como Moodle ou Canvas. Esta integração possibilita a criação de instrumentos de avaliação dinâmicos, onde quizzes automáticos buscam dados reais persistidos no servidor central para validar as respostas dos alunos, fechando o ciclo entre a experimentação física e a formalização acadêmica dentro de um fluxo de trabalho orquestrado e seguro \cite{lms}.

\subsection{Estudo de Caso: Telemetria no Circuito RC}
O circuito composto por um resistor e um capacitor, comumente denominado Circuito RC, constitui um dos pilares experimentais do ensino de eletromagnetismo e instrumentação eletrônica \cite{aguilar, gprime}. Sua relevância pedagógica reside na capacidade de ilustrar fenômenos dinâmicos de armazenamento e dissipação de energia, permitindo que o estudante observe a transição entre estados estacionários através de comportamentos temporais previsíveis. No entanto, a análise deste circuito em ambientes laboratoriais modernos transcende a mera aplicação de fórmulas, configurando-se como um problema rico para a construção de novos conhecimentos e para a validação de modelos de aquisição de dados digitais.

Sob a perspectiva da modelagem física, o processo de carga e descarga de um capacitor é regido pela segunda lei de Kirchhoff (lei das malhas), que resulta em uma equação diferencial ordinária (EDO) de primeira ordem. O modelo teórico tradicional prevê que a tensão no capacitor varie de forma exponencial em função do tempo, definida pela constante de tempo, que representa o intervalo necessário para que a carga atinja aproximadamente 63,2\% de seu valor final. Contudo, trabalhos recentes como o de Kowalski (2024) destacam que o uso de instrumentação de alta resolução revela discrepâncias significativas entre os dados experimentais e os modelos de livros didáticos. Variáveis como a capacitância intrínseca de interruptores, a autoindutância de loops de fios e a resistência em série do próprio capacitor (ESR) introduzem comportamentos não ideais que desafiam o estudante a reconciliar a teoria com a realidade física observada.

A transição da observação manual para a digitalização do fenômeno é o ponto de convergência com a área de Sistemas de Informação. Tradicionalmente, estudantes realizavam medições pontuais com multímetros e cronômetros manuais, um processo sujeito a erros humanos e com baixa densidade amostral, e portanto, baixo feedback de resultados. O surgimento de técnicas de análise de vídeo, como o uso do software Tracker, permitiu a automação da coleta de dados a partir da deflexão de ponteiros analógicos, aumentando a quantidade de pontos experimentais de acordo com a taxa de quadros (fps) da câmera utilizada \cite{aguilar}. Entretanto, a telemetria baseada em microcontroladores representa um salto qualitativo, ao realizar a conversão analógico-digital (ADC) diretamente no ponto de origem do sinal.
